% GENERATED FILE. Edit canonical lessons, then run npm run generate:books.
% canonical-source-hash: fnv1a64:723523c3a24d2015
% canonical-lessons: TA-C355-parinturai, TA-C355-arivurai, TA-C355-arivuruttu, TA-C355-aloci, TA-C355-munmoli

\chapter{Advice and Proposals}
\label{ch:ta-advice-and-proposals}
\hlchaptermodality{355}

\begin{chapteropening}
\textbf{By the end of this chapter:} \emph{I can recommend, advise and propose.}

Complete the last lesson of chapter 355: I can recommend, advise and propose.
\end{chapteropening}

\section[\texorpdfstring{\ta{பரிந்துரை} (parinturai)}{பரிந்துரை (parinturai)}]{\ta{பரிந்துரை} (parinturai) — to recommend}

\label{lesson:TA-C355-parinturai}



\begin{quote}
Before the new one: say the Tamil for a need, then the Tamil for to stress.
\end{quote}

\subsection*{You'll want to know: \ta{பரிந்துரை}}
\textbf{\ta{பரிந்துரை}} — \emph{parinturai} — \textquotedblleft{}to recommend\textquotedblright{}.

\begin{quote}
\ta{மருத்துவர்} \ta{இந்த} \ta{மருந்தைப்} \ta{பரிந்துரைத்தார்}.
\end{quote}

\begin{quote}
\emph{maruttuvar inta maruntaip parinturaittār.}
\end{quote}

\begin{quote}
The doctor recommended this medicine.
\end{quote}

Past (I …): \ta{பரிந்துரைத்தேன்}, \emph{parinturaittēṉ}.

The same word is also a noun, a recommendation: \ta{அவருடைய} \ta{பரிந்துரை}.

\begin{cousinweb}
\textbf{\ta{பரிந்துரை}} (\emph{parinturai}) is native: \emph{parintu}, \textquotedblleft{}having felt for\textquotedblright{}, from \emph{pari}, \textquotedblleft{}to be affectionate\textquotedblright{} (DEDR 3964), joined to \emph{urai}, \textquotedblleft{}to speak\textquotedblright{} (DEDR 648): to speak up for someone or something.
\end{cousinweb}

\begin{grammarlens}[title={What you've built}]
One more word for giving reasons and opinions.
\end{grammarlens}

\subsection*{Your turn}
\begin{itemize}
\raggedright
  \item \emph{Say it:} \emph{parinturai}
  \item \emph{Say it:} \emph{parinturai}, once more
  \item \emph{Say it:} \emph{maruttuvar inta maruntaip parinturaittār.} — The doctor recommended this medicine.
  \item \emph{Recall:} say the Tamil for a need, then the Tamil for to stress, then say \emph{parinturai} again
\end{itemize}

\begin{tcolorbox}[breakable,colback=teal!4,colframe=teal!35!black,title={Before you move on}]
What does \ta{பரிந்துரை} mean? (\textquotedblleft{}To recommend\textquotedblright{}.) Say it once more.
\end{tcolorbox}
\section[\texorpdfstring{\ta{அறிவுரை} (aṟivurai)}{அறிவுரை (aṟivurai)}]{\ta{அறிவுரை} (aṟivurai) — advice}

\label{lesson:TA-C355-arivurai}



\begin{quote}
Before the new one: say the Tamil for to stress, then the Tamil for to recommend.
\end{quote}

\subsection*{You'll want to know: \ta{அறிவுரை}}
\textbf{\ta{அறிவுரை}} — \emph{aṟivurai} — \textquotedblleft{}advice\textquotedblright{}.

\begin{quote}
\ta{என்} \ta{அப்பா} \ta{எனக்கு} \ta{நல்ல} \ta{அறிவுரை} \ta{சொன்னார்}.
\end{quote}

\begin{quote}
\emph{eṉ appā eṉakku nalla aṟivurai coṉṉār.}
\end{quote}

\begin{quote}
My father gave me good advice.
\end{quote}

X-\ta{க்கு} \ta{அறிவுரை} \ta{சொல்} or \ta{கூறு}, to advise X.

\begin{cousinweb}
\textbf{\ta{அறிவுரை}} (\emph{aṟivurai}) is native: \emph{aṟivu}, \textquotedblleft{}knowledge, wisdom\textquotedblright{} (DEDR 314), and \emph{urai}, \textquotedblleft{}word, utterance\textquotedblright{} (DEDR 648): words of wisdom.
\end{cousinweb}

\begin{grammarlens}[title={What you've built}]
One more word for giving reasons and opinions.
\end{grammarlens}

\subsection*{Your turn}
\begin{itemize}
\raggedright
  \item \emph{Say it:} \emph{aṟivurai}
  \item \emph{Say it:} \emph{aṟivurai}, once more
  \item \emph{Say it:} \emph{eṉ appā eṉakku nalla aṟivurai coṉṉār.} — My father gave me good advice.
  \item \emph{Recall:} say the Tamil for to stress, then the Tamil for to recommend, then say \emph{aṟivurai} again
\end{itemize}

\begin{tcolorbox}[breakable,colback=teal!4,colframe=teal!35!black,title={Before you move on}]
What does \ta{அறிவுரை} mean? (\textquotedblleft{}Advice\textquotedblright{}.) Say it once more.
\end{tcolorbox}
\section[\texorpdfstring{\ta{அறிவுறுத்து} (aṟivuṟuttu)}{அறிவுறுத்து (aṟivuṟuttu)}]{\ta{அறிவுறுத்து} (aṟivuṟuttu) — to advise; to instruct}

\label{lesson:TA-C355-arivuruttu}



\begin{quote}
Before the new one: say the Tamil for to recommend, then the Tamil for advice.
\end{quote}

\subsection*{You'll want to know: \ta{அறிவுறுத்து}}
\textbf{\ta{அறிவுறுத்து}} — \emph{aṟivuṟuttu} — \textquotedblleft{}to advise; to instruct\textquotedblright{}.

\begin{quote}
\ta{காவலர்கள்} \ta{மக்களை} \ta{வீட்டிலேயே} \ta{இருக்குமாறு} \ta{அறிவுறுத்தினார்கள்}.
\end{quote}

\begin{quote}
\emph{kāvalarkaḷ makkaḷai vīṭṭilēyē irukkumāṟu aṟivuṟuttiṉārkaḷ.}
\end{quote}

\begin{quote}
The police advised people to stay at home.
\end{quote}

Past (I …): \ta{அறிவுறுத்தினேன்}, \emph{aṟivuṟuttiṉēṉ}.

The person takes -\ta{ஐ} and the action -\ta{உமாறு}: \ta{இருக்குமாறு} \ta{அறிவுறுத்து}, to advise to stay.

\begin{cousinweb}
\textbf{\ta{அறிவுறுத்து}} (\emph{aṟivuṟuttu}) is native: \emph{aṟivu}, \textquotedblleft{}knowledge\textquotedblright{} (DEDR 314), and \emph{uṟuttu}, \textquotedblleft{}to impress strongly upon the mind\textquotedblright{} (DEDR 712): to press knowledge on someone.
\end{cousinweb}

\begin{grammarlens}[title={What you've built}]
One more word for giving reasons and opinions.
\end{grammarlens}

\subsection*{Your turn}
\begin{itemize}
\raggedright
  \item \emph{Say it:} \emph{aṟivuṟuttu}
  \item \emph{Say it:} \emph{aṟivuṟuttu}, once more
  \item \emph{Say it:} \emph{kāvalarkaḷ makkaḷai vīṭṭilēyē irukkumāṟu aṟivuṟuttiṉārkaḷ.} — The police advised people to stay at home.
  \item \emph{Recall:} say the Tamil for to recommend, then the Tamil for advice, then say \emph{aṟivuṟuttu} again
\end{itemize}

\begin{tcolorbox}[breakable,colback=teal!4,colframe=teal!35!black,title={Before you move on}]
What does \ta{அறிவுறுத்து} mean? (\textquotedblleft{}To advise; to instruct\textquotedblright{}.) Say it once more.
\end{tcolorbox}
\section[\texorpdfstring{\ta{ஆலோசி} (ālōci)}{ஆலோசி (ālōci)}]{\ta{ஆலோசி} (ālōci) — to consult; to deliberate}

\label{lesson:TA-C355-aloci}



\begin{quote}
Before the new one: say the Tamil for advice, then the Tamil for to advise.
\end{quote}

\subsection*{You'll want to know: \ta{ஆலோசி}}
\textbf{\ta{ஆலோசி}} — \emph{ālōci} — \textquotedblleft{}to consult; to deliberate\textquotedblright{}.

\begin{quote}
\ta{முடிவு} \ta{எடுப்பதற்கு} \ta{முன்} \ta{குடும்பத்துடன்} \ta{ஆலோசித்தேன்}.
\end{quote}

\begin{quote}
\emph{muṭivu eṭuppataṟku muṉ kuṭumpattuṭaṉ ālōcittēṉ.}
\end{quote}

\begin{quote}
I consulted my family before making a decision.
\end{quote}

Past (I …): \ta{ஆலோசித்தேன்}, \emph{ālōcittēṉ}.

Who is consulted takes -\ta{உடன்} or -\ta{இடம்}: \ta{மருத்துவரிடம்} \ta{ஆலோசி}.

\begin{cousinweb}
\textbf{\ta{ஆலோசி}} (\emph{ālōci}) is a verb built on Sanskrit \emph{ālocana}, \textquotedblleft{}seeing; considering, reflection\textquotedblright{} (Monier-Williams), with the loan-verb ending -\ta{இ}. The noun \ta{ஆலோசனை} means \textquotedblleft{}advice\textquotedblright{}.
\end{cousinweb}

\begin{grammarlens}[title={What you've built}]
One more word for giving reasons and opinions.
\end{grammarlens}

\subsection*{Your turn}
\begin{itemize}
\raggedright
  \item \emph{Say it:} \emph{ālōci}
  \item \emph{Say it:} \emph{ālōci}, once more
  \item \emph{Say it:} \emph{muṭivu eṭuppataṟku muṉ kuṭumpattuṭaṉ ālōcittēṉ.} — I consulted my family before making a decision.
  \item \emph{Recall:} say the Tamil for advice, then the Tamil for to advise, then say \emph{ālōci} again
\end{itemize}

\begin{tcolorbox}[breakable,colback=teal!4,colframe=teal!35!black,title={Before you move on}]
What does \ta{ஆலோசி} mean? (\textquotedblleft{}To consult; to deliberate\textquotedblright{}.) Say it once more.
\end{tcolorbox}
\section[\texorpdfstring{\ta{முன்மொழி} (muṉmoḻi)}{முன்மொழி (muṉmoḻi)}]{\ta{முன்மொழி} (muṉmoḻi) — to propose}

\label{lesson:TA-C355-munmoli}



\begin{quote}
Before the new one: say the Tamil for to advise, then the Tamil for to consult.
\end{quote}

\subsection*{You'll want to know: \ta{முன்மொழி}}
\textbf{\ta{முன்மொழி}} — \emph{muṉmoḻi} — \textquotedblleft{}to propose\textquotedblright{}.

\begin{quote}
\ta{அவர்} \ta{ஒரு} \ta{புதிய} \ta{யோசனையை} \ta{முன்மொழிந்தார்}.
\end{quote}

\begin{quote}
\emph{avar oru putiya yōcaṉaiyai muṉmoḻintār.}
\end{quote}

\begin{quote}
He proposed a new idea.
\end{quote}

Past (I …): \ta{முன்மொழிந்தேன்}, \emph{muṉmoḻintēṉ}.

What is proposed takes -\ta{ஐ}. In a meeting, one person proposes and another seconds.

\begin{cousinweb}
\textbf{\ta{முன்மொழி}} (\emph{muṉmoḻi}) is native: \emph{muṉ}, \textquotedblleft{}in front, before\textquotedblright{} (DEDR 5020), and \emph{moḻi}, \textquotedblleft{}to say, speak\textquotedblright{} (DEDR 4989): to say something first, before the others.
\end{cousinweb}

\begin{grammarlens}[title={What you've built}]
One more word for giving reasons and opinions.
\end{grammarlens}

\subsection*{Your turn}
\begin{itemize}
\raggedright
  \item \emph{Say it:} \emph{muṉmoḻi}
  \item \emph{Say it:} \emph{muṉmoḻi}, once more
  \item \emph{Say it:} \emph{avar oru putiya yōcaṉaiyai muṉmoḻintār.} — He proposed a new idea.
  \item \emph{Recall:} say the Tamil for to advise, then the Tamil for to consult, then say \emph{muṉmoḻi} again
\end{itemize}

\begin{tcolorbox}[breakable,colback=teal!4,colframe=teal!35!black,title={Before you move on}]
What does \ta{முன்மொழி} mean? (\textquotedblleft{}To propose\textquotedblright{}.) Say it once more.
\end{tcolorbox}
